\begin{abstract}
\lisa{"abstract" vs "portable" - on second thought, if you define abstract, it's fine to leave it :)}
\natalie{What is the definition of "intent"? It is loaded. I'm concerned that "intent" makes a stronger psychological claim than our evidence supports. It implies not only a plan pursuing a goal but also hints at a subjective experience. Consider, for example, a billiard ball moving toward another ball following the laws of physics. Do you say that the ball *intends* to strike the other ball, or do you say that it is *going* to strike it? In this sense, when we relate intentions for LLMs, we are overclaiming, as we have no way—nor is any known—to measure that subjective experience. My proposal: replace the word "intent" with "plan" or "prediction".}
When a Large Language Model (LLM) writes code to satisfy a multi-step specification, does it know \emph{what it intends to do next} before it starts writing? \lisa{suggestion: does it plan each step before outputting/generating an implementation} We show that code LLMs encode a compact representation of \lisa{could drop "a compact representation of" as it's already covered by "encode"} their next instrumental sub-goal before emitting its implementation, and that this encoding is \emph{decodable} and \emph{portable}. Using causal mediation analysis, we localize this encoding to a set of attention heads, which we call the \goalheads. We introduce \goalscope, a lightweight probe that uses what these heads write to the \lisa{alt suggestion: We introduce \goalscope, a lightweight probe trained on the residual stream to ... } residual stream to verbalize the upcoming sub-task in natural language. We demonstrate that the encoding is abstract: it can be extracted from one context and transplanted to trigger the same sub-goal in a different context and programming language. These heads are a crucial bottleneck in how the LM communicates the intent encoding to later layers for its concrete realization: preserving their access to task information largely preserves the capability of executing the next sub-task, while restricting that access in a broader set of goal heads sharply impairs it. We replicate the readout and transfer findings across LLM families and sizes, and explore the same heads' readouts in non-code settings. Our results provide concrete, testable evidence that LLMs encode abstract representations of the next instrumental goal before acting on it.
\lisa{minor comment: i'd stay consistent between embedding and representation.}
\lisa{looks great!}
\end{abstract}
